\emph{Preliminaries.} The properties of the adjoint and of the scalar part used are those of the preliminaries of Theorem \claimref{synthesis.chsh-bound-for-two-subsystems-of-a-clifford-algebra-by-a-rational-certificate}: $\langle xy\rangle_{0} = \langle yx\rangle_{0}$, $\langle x^{\dagger}\rangle_{0} = \langle x\rangle_{0}$, $\langle xx^{\dagger}\rangle_{0}$ is the sum of the squares of the coefficients of $x$ (so it vanishes only for $x = 0$), and the scalar part of a product of two blades is zero unless they are equal up to sign; the blades are a basis (Theorem \claimref{synthesis.signed-directions-and-the-integer-clifford-algebra}). The $p_{j}$ are commuting self-adjoint elements of square $1$; $p_{j}$ anticommutes with $e_{j}$, $f_{j}$ and commutes with the other generators, so $\Gamma$ is self-adjoint of square $1$ and anticommutes with every generator: $\alpha$ is the grade involution, an automorphism commuting with $\dagger$, and $E$ is the projection onto the even part, the group average of Theorem \claimref{singularity.descent-parity-average} for the group $\{\mathrm{id}, \alpha\}$. Write $x_{+} = E(x)$ and $x_{-} = x - E(x)$. For a self-adjoint blade $s$, $s^{2} = 1$ and $1 \pm s = \tfrac12(1 \pm s)(1 \pm s)^{\dagger}$, so $x \mapsto \langle(1 \pm s)x\rangle_{0}$ satisfies (P) (as does $\langle x\rangle_{0}$), and it is even when $s$ is even.

\emph{(i).} $P_{\pm}$ are self-adjoint idempotents with $P_{+}P_{-} = 0$, $P_{+} + P_{-} = 1$, $a = P_{+} - P_{-}$; expanding gives $P_{+}xP_{+} + P_{-}xP_{-} = \tfrac12(x + axa)$. The listed properties follow: idempotence from $P_{\pm}P_{\mp} = 0$; $E_{a}(xx^{\dagger}) = \sum_{\pm}P_{\pm}xx^{\dagger}P_{\pm} = \sum_{\pm}(P_{\pm}x)(P_{\pm}x)^{\dagger}$. If $\omega$ satisfies (P) then $\omega(E_{a}(xx^{\dagger})) \ge 0$, $\omega(E_{a}(1)) = 1$ and $\omega(E_{a}(x^{\dagger})) = \omega(E_{a}(x)^{\dagger}) = \omega(E_{a}(x))$.

\emph{(ii).} (3) implies (2): $\alpha(E_{a}(x)) = \tfrac12(\alpha(x) + \alpha(a)\alpha(x)\alpha(a)) = E_{a}(\alpha(x))$ when $\alpha(a) = \pm a$. (2) implies (1): $\omega \circ E_{a} \circ \alpha = \omega \circ \alpha \circ E_{a} = \omega \circ E_{a}$, and $\omega \circ E_{a}$ is a state by (i). (1) implies (3): write $a = a_{+} + a_{-}$; both parts are self-adjoint, and $a^{2} = 1$ gives $a_{+}^{2} + a_{-}^{2} = 1$ and $a_{+}a_{-} + a_{-}a_{+} = 0$. \emph{Step 1.} Let $y$ be odd and self-adjoint and $w = E(E_{a}(y))$, even and self-adjoint. For every even state $\omega$, $\omega(E_{a}(y)) = 0$ by (1), and $\omega(E_{a}(y)) = \omega(w)$ since $\omega$ is even. Taking $\omega = \langle\,\cdot\,\rangle_{0}$ and $\omega = \langle(1 \pm s)\,\cdot\,\rangle_{0}$ for the even self-adjoint blades $s$ gives $\langle w\rangle_{0} = 0$ and $\langle sw\rangle_{0} = 0$; a self-adjoint even element is a combination of the scalar and of even self-adjoint blades (a blade is self-adjoint or anti-self-adjoint), so $w = 0$. \emph{Step 2.} Let $x$ be even and self-adjoint and $z = (E_{a}(x))_{-} = \tfrac12(a_{+}xa_{-} + a_{-}xa_{+})$, odd and self-adjoint. For odd self-adjoint $y$, $\langle zy\rangle_{0} = \langle E_{a}(x)y\rangle_{0} = \langle xE_{a}(y)\rangle_{0} = \langle xw\rangle_{0} = 0$ (the scalar part of an even times an odd element is zero; $\langle E_{a}(x)y\rangle_{0} = \langle xE_{a}(y)\rangle_{0}$ by cyclicity; Step 1). For odd anti-self-adjoint $y$, $\langle zy\rangle_{0} = \langle(zy)^{\dagger}\rangle_{0} = -\langle yz\rangle_{0} = -\langle zy\rangle_{0}$, so it is zero. Hence $\langle zz^{\dagger}\rangle_{0} = 0$ and $z = 0$:
\[
a_{+}xa_{-} + a_{-}xa_{+} = 0 \qquad \text{for every even self-adjoint } x . \tag{$*$}
\]
\emph{Step 3.} Put $x = a_{+}$ in $(*)$: $a_{+}^{2}a_{-} + a_{-}a_{+}^{2} = 2a_{-} - 2a_{-}^{3} = 0$, so $a_{-}^{3} = a_{-}$ and $q = a_{-}^{2}$ is a self-adjoint even idempotent commuting with $a_{-}$; it commutes with $a_{+}$ (two anticommutations), and $(a_{+}q)(a_{+}q)^{\dagger} = a_{+}^{2}q = (1 - q)q = 0$, so $a_{+}q = 0$. Thus $a_{+} = a_{+}(1 - q)$, $a_{-} = a_{-}q$, $a_{+}^{2} = 1 - q$, $a_{-}^{2} = q$. \emph{Step 4.} For even $x'$ the element $x = qx'(1 - q) + (1 - q)x'^{\dagger}q$ is even and self-adjoint, and $(*)$ reads $a_{+}(1 - q)x'^{\dagger}q\,a_{-} + a_{-}q\,x'(1 - q)a_{+} = 0$. The first term lies in $(1 - q)\mathcal{A}q$ and the second in $q\mathcal{A}(1 - q)$ (with $\mathcal{A}$ the algebra), so each is zero; multiplying $a_{-}qx'(1 - q)a_{+} = 0$ by $a_{-}$ on the left and $a_{+}$ on the right gives $qx'(1 - q) = 0$, and the adjoint statement gives $(1 - q)x'q = 0$. So $q$ commutes with every even element. \emph{Step 5.} An even element commuting with every even element commutes with every product $g_{i}g_{j}$ of two distinct generators; conjugation by $g_{i}g_{j}$ multiplies a blade by $-1$ exactly when the blade contains one of the two generators and not the other, so only the scalar blade and the top blade survive: $q = c_{0} + c_{1}\Gamma$, and $q^{2} = q$ gives $q \in \{0, 1, \tfrac12(1 \pm \Gamma)\}$. If $q = \tfrac12(1 \pm \Gamma)$ then $a_{-}q = (1 - q)a_{-}$ because $a_{-}$ anticommutes with $\Gamma$, while $a_{-}q = qa_{-}$; multiplying by $q$ gives $qa_{-} = 0$, so $a_{-} = a_{-}^{3} = a_{-}q = qa_{-} = 0$ and then $q = a_{-}^{2} = 0$, a contradiction. So $q = 0$, and $a_{-}a_{-}^{\dagger} = q = 0$ gives $a = a_{+}$; or $q = 1$, and $a_{+}a_{+}^{\dagger} = 1 - q = 0$ gives $a = a_{-}$.

\emph{(iii).} The self-adjoint blades of the algebra of mode $1$ are $1$, $e_{1}$, $e_{1}f_{1} = -p_{1}$ ($f_{1}$ is anti-self-adjoint), so its self-adjoint elements are $t + xe_{1} + zp_{1}$; the square is $t^{2} + x^{2} + z^{2} + 2t(xe_{1} + zp_{1})$, which is $1$ only for $\pm 1$ or $t = 0$, $x^{2} + z^{2} = 1$. The even part is $zp_{1}$ and the odd part $xe_{1}$, so the homogeneous ones are $\pm e_{1}$, $\pm p_{1}$. The twisted mode is the image of mode $2$ under the isomorphism $e_{2} \mapsto e_{2}p_{1}$, $f_{2} \mapsto f_{2}p_{1}$ of Theorem \claimref{tower.commuting-copies}, which preserves $\dagger$, parity and $p_{2}$; so its homogeneous observables are $\pm e_{2}p_{1}$, $\pm p_{2}$. A product $A_{i}B_{j}$ is even exactly when both factors are odd (the product is $\pm U$) or both even ($\pm V$), so $E(C) = \mu U + \nu V$; the check enumerates the $256$ choices and finds $|\mu| + |\nu| \le 2$ in each. $U$, $V$ are commuting self-adjoint elements of square $1$, and for signs $\sigma, \tau$
\[
2 - \sigma U - \tau V = \tfrac12(1 - \sigma U)(1 - \sigma U)^{\dagger} + \tfrac12(1 - \tau V)(1 - \tau V)^{\dagger} ,
\]
so a state has $\omega(\sigma U) \le 1$, $\omega(\tau V) \le 1$ for all signs and $|\omega(\mu U + \nu V)| \le |\mu| + |\nu| \le 2$. For an even $\omega$, $\omega(C) = \omega(E(C))$.

\emph{(iv).} The state, the average $E(C) = \tfrac75(U + V)$ and the value $\tfrac{14}{5}$ are those of clause (v) of Theorem \claimref{synthesis.positivity-of-values-is-additivity-of-non-negative-weights-on-effects} at these directions, with $u = e_{2}f_{1} = U$ and $\Gamma = V$ (the check recomputes the value). $B_{1}$ has the odd part $\tfrac35X_{2}$ and the even part $\tfrac45Z_{2}$, so by (ii) some even state is sent to a functional that is not even; for this $\omega$ the value $\tfrac{12}{25}$ of $\omega \circ E_{B_{1}}$ on $X_{1}Z_{2}$ is computed by the check.

\emph{(v).} $p_{1}p_{2}$ is even, self-adjoint, of square $1$, anticommutes with $e_{1}$, $e_{2}$ and commutes with $e_{3}$, $e_{4}$; so $e_{3}p_{1}p_{2}$ and $e_{4}p_{1}p_{2}$ are odd, self-adjoint, of square $1$, anticommute with each other and commute with $e_{1}$ and $e_{2}$ (two anticommutations each), by the twist of Theorem \claimref{tower.commuting-copies} applied to the two modes of party $1$ together. The configuration is one of blades with $A_{1}A_{2} = -A_{2}A_{1}$, $B_{1}B_{2} = -B_{2}B_{1}$, so clause (iv) of the CHSH theorem applies: $C^{2} = 4 + 4K$, $KC = C$, and $\langle(1 + K)(1 + rC)x\rangle_{0}$ satisfies (P) for $8r^{2} \le 1$ with value $8r$ on $C$; at $r = \tfrac13$ the check also verifies $(1 + K)(1 + \tfrac13C) = \tfrac34R^{2}$ with $R = (1 + K)(\tfrac23 + \tfrac16C)$ self-adjoint. $C$ and $K$ are products of an even number of odd elements, so the functional is even. The observables are odd, so their record maps commute with $\alpha$ by (ii).
